\chapter{Station Strategie}

\section{Mandat}
La station Strategie transforme les findings de la station Signaux en un plan
tradable complet : univers, regles d'entree/sortie, contraintes de risque,
parametres optimisables et bornes de recherche.

\section{Transformation signal vers regle}
La logique est de mapper un score continu vers une action de portefeuille. Une
forme operationnelle utilisee est :
\begin{equation}
\text{position\_target}_t = \mathrm{clip}\left(\beta \cdot
\text{score}_t,\,-1,\,1\right)
\label{eq:position-target}
\end{equation}
avec $\beta$ reglant la sensibilite. L'equation~\ref{eq:position-target}
structure la transition entre information et execution.

\section{Sections de configuration}
Le Tableau~\ref{tab:strategy-sections} presente les six sections de
construction d'une strategie.

\begin{table}[H]
\centering
\caption{Sections de la station Strategie}
\label{tab:strategy-sections}
\begin{tabular}{p{3.8cm}p{8.8cm}}
\toprule
\textbf{Section} & \textbf{Decision prise} \\
\midrule
Type de strategie & Cadre logique (tendance, mean reversion, hybride). \\
Signal construction & Familles et parametres utilises pour produire le score. \\
Regles d'entree & Conditions d'ouverture et taille initiale. \\
Regles de sortie & Conditions de reduction/fermeture de position. \\
Risque & Stop loss, take profit, cooldown, limites d'exposition. \\
Review \& handoff & Verification de coherence et export WFO-ready. \\
\bottomrule
\end{tabular}
\end{table}

Le Tableau~\ref{tab:strategy-sections} est la grammaire de configuration
utilisee avant tout lancement de backtest.

\section{Dimensionnement et risque}
Le dimensionnement combine la conviction signal, les contraintes de liquidite et
les limites portefeuille. Une formulation simple de taille est :
\begin{equation}
\text{size}_t = \min\left(\text{size}_{\max},
\gamma\cdot |\text{score}_t| \cdot \text{capital}_t\right)
\label{eq:sizing}
\end{equation}
ou $\gamma$ est un facteur de calibration. L'equation~\ref{eq:sizing}
garantit une croissance graduelle de l'exposition.

\section{Handoff vers Backtest}
La station Strategie remet un objet de handoff qui contient :
\begin{itemize}
\item les regles definitives par action,
\item la liste des parametres marques optimisables,
\item les bornes et pas de recherche associes,
\item les avertissements de validite (donnees insuffisantes, complexite excessive).
\end{itemize}
Cette sortie constitue l'entree contractuelle du moteur WFO.

